\chapter{How Research Works}
\label{ch:research}

\begin{center}\itshape
Choose $n$ positive integers so that no two different selections of them have the same sum.\\
How small can the largest one be? A young Erd\H os asked in 1931.\\
Ninety-five years later the answer is still not known, and the story is not over.
\end{center}

Every research problem has a life. Somebody asks a question, often in passing. A first argument gives a bound; somebody else
finds a construction going the other way; the two bounds refuse to meet; the question is written down in a list of open
problems, a prize is attached, and for decades people chip away at the constants. Then, one day, something new happens. This
chapter follows one problem through that whole life, because it is the best way to see what researchers actually do all day:
they read, they compute, they bound, they construct, they check, and they write.

The problem is the one Paul Erd\H os called ``perhaps my first serious problem''. In a paper of 1998 he dated it to
1931, when he was eighteen; the quotation and the date are recorded on the first page of the Erd\H os problems database, whose
problem number~$1$ it is. Erd\H os offered $500$ dollars for its solution. His first result on it, with Leo Moser, was announced at
a colloquium in Brussels in 1955 and printed in 1956; in 1985 he offered another $100$ dollars to anyone who improved the constant
$\frac14$ in their bound. The constant was improved by Guy, Elkies, Bae, Aliev, Elkies and Gleason, and finally Dubroff, Fox and
Xu, whose two-page note of 2021 is the paper we shall read line by line in the next chapter.

And in September 2026, as this book was being written, the conjecture fell, disproved by a construction found by an
artificial intelligence system and explained to humans by the mathematician who runs the database. We shall check some pieces of
that construction ourselves at the end of the chapter, because checking a new proof is also part of the job.

\begin{reminder}[from Volume~4 and its combinatorial appendices]
\emph{Chebyshev's inequality} (Volume~4, Chapter~4): $P(|X-\mathbb EX|\ge t)\le\operatorname{Var}X/t^2$. \emph{Variance of a sum of independent
variables} is the sum of the variances. A \emph{lattice} in $\R^d$ is a discrete subgroup; its \emph{covolume} is the volume of a fundamental domain,
$|\det M|$ when it is $M\Z^d$ (Volume~13, Chapter~5).
\end{reminder}

\section{The problem}

\begin{definition}\label{def:dissociated}
A finite set $A$ of positive integers is \emph{dissociated}, or has \emph{distinct subset sums}, if the $2^{|A|}$ sums
$\sum_{a\in S}a$, $S\subseteq A$, are pairwise different. For $n\ge1$ let $a(n)$ be the least possible value of $\max A$ over all dissociated sets $A$
with $n$ elements.
\end{definition}

The powers of $2$, $\{1,2,4,\dots,2^{n-1}\}$, are dissociated: a subset sum is a number written in binary, and binary
expansions are unique. So $a(n)\le2^{n-1}$. Erd\H os's conjecture was that this is the truth up to a constant:
\[
a(n)\ge c\,2^n\quad\text{for some absolute constant }c>0.
\]

An equivalent way to say ``dissociated'' will be used all the time: $A=\{a_1,\dots,a_n\}$ is dissociated if and only if
\begin{equation}\label{eq:dissocsigned}
\sum_{i=1}^n c_ia_i\neq0\quad\text{for every }c\in\{-1,0,1\}^n\setminus\{0\}.
\end{equation}
Indeed two subsets $S\ne T$ have equal sums exactly when $\sum_{S\setminus T}a_i-\sum_{T\setminus S}a_i=0$, which is a relation~\eqref{eq:dissocsigned} with $c_i=1$ on $S\setminus T$, $-1$ on
$T\setminus S$ and $0$ elsewhere; and every nonzero $c\in\{-1,0,1\}^n$ arises in this way from $S=\{c_i=1\}$, $T=\{c_i=-1\}$.

\begin{think}
Before reading on: can you beat the powers of $2$ for $n=4$, that is, find four positive integers, all at most $7$, with distinct
subset sums? And why can't you do it with four integers at most $5$?
\end{think}

\begin{proposition}[the counting bound]\label{prop:counting}
$a(n)\ge\dfrac{2^n-1}n$.
\end{proposition}

\begin{proof}
If $A$ is dissociated with $n$ elements and largest element $N$, its $2^n$ subset sums are distinct integers in $\{0,1,\dots,nN\}$, a set of
$nN+1$ elements. So $2^n\le nN+1$.
\end{proof}

So $\frac{2^n}n\lesssim a(n)\le2^{n-1}$, and the whole problem lives in the factor $n$ between the two bounds.

\begin{proposition}[doubling]\label{prop:doubling}
If $A$ is dissociated, so is $\{1\}\cup2A$. Hence $a(n+1)\le2a(n)$.
\end{proposition}

\begin{proof}
A subset sum of $\{1\}\cup2A$ is $\varepsilon+2s$ with $\varepsilon\in\{0,1\}$ and $s$ a subset sum of $A$. Its parity gives $\varepsilon$, then $s$ is determined, and $s$
determines the subset of $A$. The largest element of $\{1\}\cup2A$ is $2\max A$ (if $\max A\ge1$).
\end{proof}

Small cases can be found by hand or by a short search. For $n=4$ the set $\{3,5,6,7\}$ works: its sixteen subset sums
\[
0,3,5,6,7,8,9,10,11,12,13,14,15,16,18,21
\]
are distinct, and no set inside $\{1,\dots,6\}$ works (Exercise~\ref{ex:afour}). Exhaustive searches give the first values:
\begin{center}
\begin{tabular}{c|c|c||c|c|c}
$n$&$a(n)$&$a(n)/2^n$&$n$&$a(n)$&$a(n)/2^n$\\\hline
$1$&$1$&$0.50$&$6$&$24$&$0.38$\\
$2$&$2$&$0.50$&$7$&$44$&$0.34$\\
$3$&$4$&$0.50$&$8$&$84$&$0.33$\\
$4$&$7$&$0.44$&$9$&$161$&$0.31$\\
$5$&$13$&$0.41$&$10$&$309$&$0.30$\\
\end{tabular}
\end{center}
Small cases were computed by Lunnon (1988) and others; the value $a(10)=309$ was settled by a computation of Paul Dyson in October 2025,
as recorded in the OEIS (sequence A276661), which lists the references. These values can be recomputed by a finite search.

\section{Constructions: the Conway--Guy sequence}

In 1968 John Conway and Richard Guy announced, in a short abstract in the \emph{Notices of the AMS}, a family of sets that seemed to beat the
powers of $2$ by a growing margin. Define $u_0=0$, $u_1=1$ and
\[
u_{m+1}=2u_m-u_{m-r},\qquad r=\operatorname{round}\bigl(\sqrt{2m}\bigr)\qquad(m\ge1),
\]
and let $A_n=\{u_n-u_{n-i}:1\le i\le n\}$. The first sequences are $\{1\}$, $\{1,2\}$, $\{2,3,4\}$, $\{3,5,6,7\}$, $\{6,9,11,12,13\}$, $\{11,17,20,22,23,24\}$,
and their largest elements $1,2,4,7,13,24,44,84,161,309$ are exactly the optimal values $a(n)$ of the table, for every $n\le10$. A computer
confirms that $A_n$ is dissociated for every $n\le24$ (a direct computation verifies $2^{24}$ sums), and the ratio $\max A_n/2^n$ keeps decreasing:
\begin{center}
\begin{tabular}{c|c}
$n$&$\max A_n/2^n$\\\hline
$12$&$0.2842$\\
$15$&$0.2688$\\
$18$&$0.2594$\\
$21$&$0.2529$\\
$24$&$0.2487$\\
\end{tabular}
\end{center}

\begin{remark}
Tom Bohman proved in 1996 that every Conway--Guy set $A_n$ is dissociated, and in 1998 he found a different construction giving
$a(n)\le0.22002\cdot2^n$ for all large $n$. Neither result is proved in this book; both are good examples of research papers whose
statements are easy and whose proofs are long case analyses. Conway and Guy also conjectured that their sets are optimal for all large
$n$; the counterexample of Section~\ref{sec:astra} disproves this.
\end{remark}

\section{The second moment bound}

Erd\H os and Moser's idea is probabilistic, and it was one of the first uses of what is now called the \emph{second moment method}.
Choose the signs at random. Let $\varepsilon_1,\dots,\varepsilon_n$ be independent and uniform in $\{-1,1\}$, and
\[
X=\varepsilon_1a_1+\dots+\varepsilon_na_n.
\]
Then $X$ takes each of its values with probability exactly $2^{-n}$, because $X=2s-\sum a_i$ where $s$ is the subset sum of
$\{a_i:\varepsilon_i=1\}$, and these $2^n$ sums are distinct. All values of $X$ have the parity of $\sum a_i$, so two of them differ by at least $2$.
And $X$ is centred, with variance $\sigma^2=\sum a_i^2\le na_n^2$.

So $X$ is a random variable with $2^n$ equally likely values, spread out by at least $2$, and with a small variance. That is a
contradiction if the variance is too small: a small variance concentrates the mass in a short window, but a short window
contains few values, each carrying only $2^{-n}$.

\begin{theorem}[Erd\H os and Moser, in a sharpened form]\label{thm:erdosmoser}
For every dissociated set $\{a_1<\dots<a_n\}$,
\[
a_n\ge\frac2{3\sqrt3}\cdot\frac{2^n}{\sqrt n}-\frac1{\sqrt{3n}}.
\]
In particular $a(n)\ge(0.384-o(1))\,2^n/\sqrt n$.
\end{theorem}

\begin{proof}
Let $\lambda>1$. By Chebyshev's inequality, $P(|X|<\lambda\sigma)\ge1-\lambda^{-2}$. The values of $X$ in the open interval $(-\lambda\sigma,\lambda\sigma)$ are spaced by at least $2$,
so there are at most $\lambda\sigma+1$ of them, each of probability $2^{-n}$. Hence
\[
1-\frac1{\lambda^2}\le(\lambda\sigma+1)2^{-n},\qquad\text{that is}\qquad\sigma\ge\frac{(1-\lambda^{-2})2^n-1}\lambda.
\]
The function $\lambda\mapsto(1-\lambda^{-2})/\lambda$ is maximal at $\lambda=\sqrt3$, where it equals $\frac2{3\sqrt3}$. With this choice, and $\sigma\le\sqrt n\,a_n$, we get the
inequality.
\end{proof}

Erd\H os and Moser stated the constant $\frac14$; the optimisation over $\lambda$ is a remark of the kind every reader of a paper should make
(Exercise~\ref{ex:chebopt} asks for the version with the exact count of values). Figure~\ref{fig:sums} shows what the proof sees.

\begin{figure}[ht]
\centering
\fig{ch01-sums}
\caption{Histogram of the $2^{14}$ values of $X=\sum\pm a_i$ for the Conway--Guy set $A_{14}$ (largest element $4484$), with the normal density of
the same variance. Chebyshev guarantees that at least two thirds of the values fall in the shaded window $|x|<\sqrt3\,\sigma$; in fact about
$92\%$ do. Since the values are at least $2$ apart, the window must be long, so $\sigma$, and with it $a_{14}$, cannot be too small.}
\label{fig:sums}
\end{figure}

\begin{think}
The histogram looks Gaussian. If $X$ were \emph{exactly} normal, a window of length $2\ell$ around $0$ would carry probability about
$\sqrt{2/\pi}\,\ell/\sigma$ for small $\ell$, much less than Chebyshev's bound suggests. What lower bound on $a_n$ would that give? (This is the idea of the paper of the
next chapter.)
\end{think}

\begin{beauty}
The argument never looks at a single subset. It replaces a question about $2^n$ sums by one number, the variance, and wins
because a variance is easy to compute and hard to fool. Erd\H os used the same move, a random object whose average is
computable, in a dozen other problems, and the method now fills whole books.
\end{beauty}

\section{The tools of the trade}

How would you have learned all of the above if you had met the problem yesterday? This is what a researcher does in the first week
with a new question.

\emph{Find the problem's name and its literature.} Search engines are bad at mathematics written in formulas, so search for words: here
``distinct subset sums'' or ``sum-distinct sets''. The two great reviewing databases list almost every research paper in mathematics since the
nineteenth century, with a short review: \emph{zbMATH Open} (free to everyone since January 2021) and the American Mathematical Society's
\emph{MathSciNet}. Their key feature is that they record who cites whom. From one good paper, go \emph{backward} through its references to the
history, and \emph{forward} through the papers citing it to the current state of the art. Lists of open problems help: Richard Guy's
\emph{Unsolved Problems in Number Theory} has this one as problem~C8, and since 2023 Thomas Bloom's database \texttt{erdosproblems.com} collects
more than a thousand problems of Erd\H os with their status, references, and a discussion forum added in August 2025.

\emph{Read the preprints.} Almost every paper now appears first on the \emph{arXiv}, the preprint server started by the physicist Paul Ginsparg in
August 1991. Its daily listings (for this problem \texttt{math.CO} and \texttt{math.NT}) are how researchers follow their field, months or years
before journal publication. A preprint is not refereed: read it as a colleague's letter, not as gospel.

\emph{Compute.} The Online Encyclopedia of Integer Sequences, which Neil Sloane began as a graduate student in 1964 and put online in 1996, turns
a list of numbers into a list of references: typing $1,2,4,7,13,24,44$ into it finds sequence A276661, with all the history above. Here the
computer also checks small cases, tests conjectured constructions and estimates constants, as in the tables of this chapter.

\emph{Ask.} \emph{MathOverflow}, started in September 2009 by Anton Geraschenko and David Zureick-Brown, then graduate students at Berkeley, is a
question-and-answer site for research-level questions; a precise question with a reference to what you already know often gets an answer
from an expert within a day. Seminars, now often online and recorded, and plain e-mail to the author of a paper (short, precise, polite)
are the other ways to ask.

\emph{Write it down.} Keep a research notebook: dated entries, the question as you currently understand it, every attempt including the failed
ones (especially those: a failed attempt recorded with the reason it failed saves a week a year later), the computations with the code that
produced them, and the references read, each with one line saying what it contains. Keep the same discipline for computational experiments.

\section{September 2026: the conjecture falls}\label{sec:astra}

Since 2023 the Erd\H os problems database has seen a flood of activity, much of it involving artificial intelligence, as told by
Konstantin Kakaes in \emph{Quanta Magazine} in August 2026: language models were first used to search the literature (and found that several
``open'' problems had been solved long ago, sometimes after an embarrassing claim of a new proof), then to find genuinely new proofs, checked by
humans and sometimes by proof assistants (Chapter~\ref{ch:lean} of this volume).

On the page of Problem~$1$, edited on 3 September 2026, the statement is marked as \emph{disproved}: according to the page, a system called GPT-6
Astra proved that for every $\epsilon>0$ there are arbitrarily large $n$ with $a(n)\le\epsilon\,2^n$. Thomas Bloom posted an exposition of the
construction, which he describes as a sketch, written in language quite far from the original write-up, and ``currently non-quantitative''. The
lower bound of order $2^n/\sqrt n$ is untouched, so the truth now lies somewhere between $2^n/\sqrt n$ and $o(2^n)$, and the problem has a new life.

What should a reader do with such news? Check the pieces. The exposition reduces the problem to three steps, and two of them are short
and elementary. We prove them now; the third is stated as a remark.

\begin{definition}
Let $Q\ge2$ be an integer. A finite sequence of positive integers $u_1,\dots,u_m$ is \emph{$Q$-dissociated} if $\sum c_iu_i\neq0$ for every
nonzero $c\in\Z^m$ with $|c_i|<Q$ for all $i$.
\end{definition}

So ``dissociated'' is ``$2$-dissociated'', by~\eqref{eq:dissocsigned}. Being $Q$-dissociated for a large $Q$ is a much stronger property, and the first step
shows that it is cheaper per element.

\begin{proposition}[from $Q$-dissociated to dissociated]\label{prop:qdissoc}
Let $Q=2^l$ and let $u_1,\dots,u_m$ be $Q$-dissociated, with largest element $Y$. Then the $lm$ numbers $2^ju_i$ ($0\le j<l$, $1\le i\le m$) form a
dissociated set with largest element $2^{l-1}Y<QY$. Consequently, since $2^{lm}=Q^m$,
\[
\frac{a(lm)}{2^{lm}}<\frac{QY}{Q^m}=\frac Y{Q^{m-1}}.
\]
\end{proposition}

\begin{proof}
Take $\varepsilon_{ij}\in\{-1,0,1\}$, not all zero, and suppose $\sum_{i,j}\varepsilon_{ij}2^ju_i=0$. Then $\sum_ic_iu_i=0$ with $c_i=\sum_{j<l}\varepsilon_{ij}2^j$, and
$|c_i|\le2^l-1<Q$. So all $c_i=0$ by $Q$-dissociation. But $c_i=0$ forces $\varepsilon_{ij}=0$ for all $j$: the lowest $j$ with $\varepsilon_{ij}\ne0$ would make $c_i$ equal to
$\pm2^j$ plus a multiple of $2^{j+1}$, which is not zero. So all $\varepsilon_{ij}=0$, a contradiction, and the set is dissociated by~\eqref{eq:dissocsigned} (its
elements are also distinct, since equal elements would give a relation with two nonzero coefficients). Its largest element is $2^{l-1}Y<QY$.
\end{proof}

So to make $a(n)/2^n$ small, it is enough to find, for some large $Q$, a $Q$-dissociated sequence with $Y/Q^{m-1}$ small. Trivially the powers
$1,Q,\dots,Q^{m-1}$ are $Q$-dissociated with $Y/Q^{m-1}=1$; the question is to beat the base-$Q$ digits by a large factor. The second step
translates this into the geometry of lattices.

\begin{proposition}[lattices give $Q$-dissociated sequences]\label{prop:latticeq}
Let $u\in\Z^m$ have positive coordinates, and let $L=\{c\in\Z^m:c\cdot u=0\}$. If $L\cap(-Q,Q)^m=\{0\}$, then $u_1,\dots,u_m$ is $Q$-dissociated.
\end{proposition}

\begin{proof}
A relation $\sum c_iu_i=0$ with $|c_i|<Q$ says $c\in L\cap(-Q,Q)^m$, so $c=0$.
\end{proof}

When $u$ is primitive (its coordinates have no common factor), $L$ is a lattice of rank $m-1$ in the hyperplane $u^\perp$, of covolume
$\|u\|_2$; so we want a lattice in a hyperplane with no nonzero point in the cube $(-Q,Q)^m$ and the smallest possible covolume. The
construction starts from a remarkable small lattice. Let $b\ge3$ be odd, index coordinates of $\R^b$ by $\Z/b\Z$, let $P$ be the cyclic shift
$(Pz)_i=z_{i-1}$, and $T=I+\frac12P$, so $(Tz)_i=z_i+\frac12z_{i-1}$.

\begin{lemma}[the odd cycle]\label{lem:oddcycle}
If $b$ is odd, $z\in\Z^b$, $\sum z_i=0$ and $|z_i+\frac12z_{i-1}|<1$ for all $i$, then $z=0$. Moreover $\det T=1+2^{-b}$.
\end{lemma}

\begin{proof}
If some $|z_i|\ge2$, take $i$ with $|z_i|$ maximal: then $|z_i+\frac12z_{i-1}|\ge|z_i|-\frac12|z_i|\ge1$, a contradiction. So $z_i\in\{-1,0,1\}$. If $z_i=1$ then
$1+\frac12z_{i-1}<1$ forces $z_{i-1}=-1$; if $z_i=-1$ then $z_{i-1}=1$ likewise. So as soon as one coordinate is nonzero, going backwards around the cycle
the coordinates are all nonzero and alternate in sign; after $b$ steps we return to $z_i$ with the sign changed $b$ times, an odd number, so
$z_i=-z_i$, a contradiction. Hence $z=0$ (the condition $\sum z_i=0$ is not even needed for this part; it places $z$ in the hyperplane where the
construction lives).

For the determinant, $P$ is diagonalisable with eigenvalues the $b$-th roots of unity $\omega$, so $\det T=\prod_{\omega^b=1}(1+\frac\omega2)$. Since
$\prod_\omega(x-\omega)=x^b-1$, taking $x=-2$ gives $\prod_\omega(-2-\omega)=(-2)^b-1$, so $\prod_\omega(2+\omega)=-((-2)^b-1)=2^b+1$ for odd $b$, and
$\det T=2^{-b}(2^b+1)=1+2^{-b}$.
\end{proof}

So $T(\Z^b)$ is a lattice of covolume barely above $1$ that meets the open unit cube only at $0$: almost as dense as $\Z^b$, and yet
avoiding the cube, thanks to an odd cycle of signs that cannot close up. Exercise~\ref{ex:evencycle} shows that the parity of $b$ is essential.

\begin{remark}[the rest of the construction]
Bloom's exposition then builds lattices in $\R^{b^s}$ by an iteration that feeds $T$ into itself $s$ times; each level divides the relevant
ratio by $\frac32$ and multiplies it by $(1+2^{-b})^{\text{(number of blocks)}}$, so choosing $s$ large and then $b$ large makes it as small as we like.
The last step turns these real lattices into lattices of the form $L$ above, with an integer primitive $u$; for this the exposition invokes a
density theorem for primitive lattices (citing work of Horesh and Karasik, 2023), which is why the result is not yet quantitative. We do not
prove these steps here, and the exposition itself is a sketch; at the time of writing the community is still digesting the argument, and
the forum of the problem lists people working on making it effective. This is research as it happens.
\end{remark}

\begin{think}
In Proposition~\ref{prop:qdissoc} the ratio of the new set is exactly $\frac{2^{l-1}Y}{2^{lm}}=\frac12\cdot\frac Y{Q^{m-1}}$. Starting from the trivial $Q$-dissociated
sequence $1,Q,\dots,Q^{m-1}$, which set do you get, and what ratio? So the whole game is to beat the base-$Q$ digits by a large factor.
\end{think}

\section{Exercises}

\begin{exercise}[$\star$]\label{ex:afour}
Show that no $4$-element subset of $\{1,\dots,6\}$ is dissociated, so $a(4)=7$. (No computer.)
\end{exercise}

\begin{exercise}[$\star$]\label{ex:conwayfive}
Check by hand that the Conway--Guy set $\{6,9,11,12,13\}$ is dissociated. Can you organise the check so that it uses fewer than the $\binom{32}2$ comparisons?
\end{exercise}

\begin{exercise}[$\star\star$]\label{ex:chebopt}
In the proof of Theorem~\ref{thm:erdosmoser}, replace Chebyshev's inequality by the exact identity $\mathbb EX^2=\sigma^2$ and the knowledge that $X$ has $2^n$ equally likely
values, symmetric about $0$, spaced by at least $2$. Show that $\sigma^2\ge\frac{4^n-1}3$, and deduce $a_n\ge\frac{2^n}{\sqrt{3n}}(1-o(1))$. (This is Moser's variance
bound, used in the next chapter.)
\end{exercise}

\begin{exercise}[$\star\star$]\label{ex:realversion}
Erd\H os and Graham also asked the question for real numbers: $0<a_1<\dots<a_n$ such that any two distinct subset sums differ by at least $1$.
Show that the counting bound and Theorem~\ref{thm:erdosmoser} still hold. Which step of Proposition~\ref{prop:doubling} needs a change?
\end{exercise}

\begin{exercise}[$\star\star$]\label{ex:evencycle}
Show that Lemma~\ref{lem:oddcycle} fails for every even $b$: find a nonzero $z\in\Z^b$ with $\sum z_i=0$ and $|z_i+\frac12z_{i-1}|<1$ for all $i$. Compute $\det T$ for even $b$.
\end{exercise}

\begin{exercise}[$\star\star$]\label{ex:qthree}
Show that $1,Q,\dots,Q^{m-1}$ is $Q$-dissociated but not $(Q+1)$-dissociated. Show that $(5,16,21)$ is not even $2$-dissociated, and that $(10,13,14)$ is
$4$-dissociated. Apply Proposition~\ref{prop:qdissoc} to $(10,13,14)$: which dissociated $6$-set do you get, and how does its largest element compare with the one
obtained from the base-$4$ digits $(1,4,16)$, and with $a(6)=24$?
\end{exercise}

\begin{exercise}[$\star\star\star$]\label{ex:greedy}
\emph{The greedy construction is bad.} Build a set by taking $a_1=1$ and, at each step, the least integer larger than the previous ones that keeps
the set dissociated. Show that this produces exactly the powers of $2$. Show that a dissociated set containing $1$ can nevertheless beat the powers of $2$:
$\{1,6,10,12,14\}$ is dissociated. What does the greedy algorithm give when started from $a_1=3$?
\end{exercise}

\begin{exercise}[$\star$]\label{ex:litsearch}
\emph{A literature search.} Using zbMATH Open or MathSciNet, find the papers that cite N.~D.~Elkies, ``An improved lower bound on the greatest
element of a sum-distinct set of fixed order'' (1986). Which of them improves the constant? Record the result in a research notebook entry, as
described in the chapter.
\end{exercise}

\section{Hints}

\begin{hint}{ex:afour}
A dissociated $4$-set has $16$ distinct sums between $0$ and the sum of its elements, so the sum of the set is at least $15$. Inside $\{1,\dots,6\}$ that leaves
very few sets.
\end{hint}

\begin{hint}{ex:conwayfive}
Two subset sums are equal if and only if a signed combination~\eqref{eq:dissocsigned} vanishes. Sort the $32$ sums.
\end{hint}

\begin{hint}{ex:chebopt}
Among all sets of $2^n$ reals, symmetric about $0$ and spaced by at least $2$, which minimises the sum of squares?
\end{hint}

\begin{hint}{ex:realversion}
Count with spacing instead of with integers. For doubling, keep $\{1\}\cup2A$ and check the spacing of its sums directly.
\end{hint}

\begin{hint}{ex:evencycle}
Alternate signs.
\end{hint}

\begin{hint}{ex:qthree}
For the powers, look at the lowest nonzero coefficient modulo $Q$. For $(10,13,14)$, take a relation $10c_1+13c_2+14c_3=0$ with $|c_i|\le3$ and reduce it
modulo $2$, then modulo $13$.
\end{hint}

\begin{hint}{ex:greedy}
By induction, after choosing $1,2,\dots,2^{k-1}$ every integer up to $2^k-1$ is a subset sum.
\end{hint}

\begin{hint}{ex:litsearch}
Search for the title; on the result page, the ``cited by'' or ``citations'' link lists the forward references.
\end{hint}

\section{Solutions}

\begin{solution}{ex:afour}
\idea The counting bound is weak in general but strong for small sets: it forces the total sum to be large, and inside $\{1,\dots,6\}$ only
sets with large elements have a large sum.
\writeup If $A$ is dissociated with $4$ elements, its $16$ subset sums are distinct integers between $0$ and $s=\sum A$, so $s\ge15$. The $4$-subsets of
$\{1,\dots,6\}$ with sum at least $15$ are $\{3,4,5,6\}$ (sum $18$), $\{2,4,5,6\}$ ($17$), $\{1,4,5,6\}$ and $\{2,3,5,6\}$ ($16$), and $\{1,3,5,6\}$, $\{2,3,4,6\}$ ($15$).
Each has a coincidence: $3+6=4+5$; $2+4=6$; $1+5=6$; $2+3=5$; $1+5=6$; $2+4=6$. So no $4$-subset of $\{1,\dots,6\}$ is dissociated, while
$\{3,5,6,7\}$ is, as its list of sums shows. Hence $a(4)=7$.
\end{solution}

\begin{solution}{ex:conwayfive}
\idea Compute the sums once, sort them, and look for repeats: one pass instead of all pairs.
\writeup The subset sums of $\{6,9,11,12,13\}$, sorted, are $0$, $6$, $9$, $11$, $12$, $13$, $15$, $17$, $18$, $19$, $20$, $21$, $22$, $23$, $24$, $25$, $26$, $27$, $28$, $29$, $30$, $31$, $32$, $33$, $34$, $36$, $38$, $39$, $40$, $42$, $45$, $51$. These are $32$ distinct numbers, so the set
is dissociated. Generating them takes $32$ additions and sorting $32\log_232=160$ comparisons. A numerical calculation does the same for $n=24$, with
$2^{24}$ sums.
\end{solution}

\begin{solution}{ex:chebopt}
\idea The spacing pushes the values of $X$ away from $0$; the most concentrated configuration is the progression $\pm1,\pm3,\dots$, whose variance can be
computed exactly.
\writeup The $2^n$ values of $X$ form a set symmetric about $0$ with an even number of elements, so $0$ is not one of them (the others would pair up). Let
$x_1<x_2<\dots<x_{2^{n-1}}$ be the positive values. As $-x_1$ is also a value, $2x_1\ge2$, so $x_1\ge1$, and by the spacing $x_k\ge2k-1$. Hence
\[
\sigma^2=\mathbb EX^2\ge2\cdot2^{-n}\sum_{k=1}^{2^{n-1}}(2k-1)^2=2^{1-n}\cdot\frac{2^{n-1}(2^n-1)(2^n+1)}3=\frac{4^n-1}3 ,
\]
using $\sum_{k\le m}(2k-1)^2=\frac{m(2m-1)(2m+1)}3$. With $\sigma^2\le na_n^2$ this gives $a_n\ge\sqrt{(4^n-1)/(3n)}$, better than Theorem~\ref{thm:erdosmoser}: the constant is
$\frac1{\sqrt3}=0.577\ldots$ instead of $0.384\ldots$. The moral: Chebyshev's inequality throws information away, and the exact structure of the problem
(equally likely, well spaced values) can be used directly.
\end{solution}

\begin{solution}{ex:realversion}
\idea Integers were only used to say that distinct sums are at least $1$ apart; that is exactly the new hypothesis.
\writeup The $2^n$ subset sums lie in $[0,na_n]$ and are $1$ apart, so $(2^n-1)\cdot1\le na_n$: the counting bound holds. In Theorem~\ref{thm:erdosmoser}, the values of $X=2s-\sum a_i$ are
$2$ apart, so at most $\lambda\sigma+1$ of them lie in $(-\lambda\sigma,\lambda\sigma)$, and the proof is unchanged. In Proposition~\ref{prop:doubling} the parity argument has no meaning for
real numbers, and is replaced by a spacing argument for the same set $\{1\}\cup2A$: its sums are $\varepsilon+2s$, two of them with the same $\varepsilon$ and different $s$ differ by
$2|s-s'|\ge2$, and two with different $\varepsilon$ differ by $|1\pm2(s-s')|$, which is $1$ if $s=s'$ and at least $2|s-s'|-1\ge1$ otherwise. So the doubling holds for real
sets too.
\end{solution}

\begin{solution}{ex:evencycle}
\idea The proof of the lemma used the odd length only at the end, to get a contradiction from alternating signs. With an even length,
alternating signs close up.
\writeup For even $b$, take $z_i=(-1)^i$. Then $\sum z_i=0$ and $z_i+\frac12z_{i-1}=(-1)^i(1-\frac12)=\pm\frac12$, of absolute value less than $1$. For the determinant,
$\prod_\omega(2+\omega)=-((-2)^b-1)=1-2^b$ for even $b$, so $\det T=2^{-b}(1-2^b)$ has absolute value $1-2^{-b}$. (Indeed $T$ has the eigenvalue $1-\frac12$ on the
alternating vector.)
\end{solution}

\begin{solution}{ex:qthree}
\idea $Q$-dissociation is the uniqueness of a representation with digits in $(-Q,Q)$; for the powers of $Q$ this is numeration in base $Q$. For a concrete
triple, congruences cut the finitely many candidate relations down to a few.
\writeup If $\sum c_iQ^{i-1}=0$ with $|c_i|<Q$ and $j$ is the least index with $c_j\ne0$, reducing modulo $Q^j$ gives $c_jQ^{j-1}\equiv0\pmod{Q^j}$, so $Q\mid c_j$, impossible
since $0<|c_j|<Q$. So the powers are $Q$-dissociated; the relation $Q\cdot1-1\cdot Q=0$ has coefficients of size at most $Q<Q+1$, so they are not
$(Q+1)$-dissociated. The triple $(5,16,21)$ has the relation $5+16-21=0$.

For $(10,13,14)$, suppose $10c_1+13c_2+14c_3=0$ with $|c_i|\le3$. Modulo $2$, $c_2$ is even. Modulo $13$, $10c_1+c_3\equiv0$, that is $c_3\equiv3c_1$; since
$|3c_1-c_3|\le12$, this forces $c_3=3c_1$, hence $|c_1|\le1$. Then $52c_1+13c_2=0$ gives $c_2=-4c_1$, so $c_1=0$ (as $|c_2|\le3$), and $c_2=c_3=0$. So the triple is $4$-dissociated.
With $l=2$, the proposition gives the dissociated set $\{10,13,14,20,26,28\}$ (a numerical calculation lists its $64$ distinct sums), with largest element $28$, against $32$
for the set $\{1,2,4,8,16,32\}$ obtained from the digits $(1,4,16)$. It is worse than the optimum $a(6)=24$, but it beats the powers of $2$ with an
object found by a search for a good $Q$-dissociated triple: the idea of the September 2026 construction in miniature.
\end{solution}

\begin{solution}{ex:greedy}
\idea After the first $k$ powers of $2$, every number up to $2^k-1$ is already a subset sum, so any new element below $2^k$ collides with one of them.
Greedy algorithms commit too early.
\writeup Induction: if the greedy set is $\{1,2,\dots,2^{k-1}\}$, its subset sums are all of $\{0,\dots,2^k-1\}$, so a new element $x<2^k$ equals an existing
sum, a collision with the subset $\{x\}$, while $x=2^k$ works. The set $\{1,6,10,12,14\}=\{1\}\cup2\{3,5,6,7\}$ is dissociated by Proposition~\ref{prop:doubling}, with maximum
$14<16$. Started from $3$, the greedy algorithm accepts $4$ and $5$ (the sums $0,3,4,5,7,8,9,12$ are distinct), then rejects $6,7,8,9$ ($6+3=4+5$, $7=3+4$, $8=3+5$,
$9=4+5$) and accepts $10$; from then on it doubles: $3,4,5,10,20,40,80,\dots$. Already for four elements its maximum $10$ is worse than the optimal
$7$. Good sets are found by search, or by structure, not by greed.
\end{solution}

\begin{solution}{ex:litsearch}
\idea Forward citation search is the fastest way from an old paper to the present.
\writeup In zbMATH Open, search for the title; the entry has a list of citing documents, and so does MathSciNet's. The list is long and changes over
time, so the answer is a method rather than a list; among the citing papers you will find Dubroff, Fox and Xu (2021), whose reference [5] is Elkies's
paper, and Aliev (2008), whose constant $\sqrt{3/(2\pi)}\approx0.69$ was the best published before 2021. A notebook entry might read: ``26 Oct.: forward search from
Elkies 1986. Aliev 2008 uses Siegel's lemma, constant $0.69$. DFX 2021: two proofs of $\sqrt{2/\pi}\approx0.80$, one via Berry--Esseen, one via Harper's
isoperimetric inequality. To read next: DFX.''
\end{solution}

\section*{Sources and further reading}
\addcontentsline{toc}{section}{Sources and further reading}

\begin{itemize}
\item \emph{The problem.} T.~F.~Bloom, Erd\H os Problem \#1, \url{https://www.erdosproblems.com/1}, accessed 26 October 2026: the statement, the
quotation ``perhaps my first serious problem'' and the date 1931 from P.~Erd\H os, ``Some of my new and almost new problems and results in combinatorial
number theory'' (1998), the prizes (500 dollars; 100 dollars for the constant, from Erd\H os 1985), the status ``disproved'' and Bloom's exposition
(last edited 3 September 2026). R.~K.~Guy, \emph{Unsolved Problems in Number Theory}, 3rd ed., Springer, 2004, problem C8.
\item \emph{Bounds and constructions.} P.~Erd\H os, ``Problems and results in additive number theory'', Colloque sur la Th\'eorie des Nombres, Bruxelles,
1955 (1956), 127--137 (the Erd\H os--Moser bound). J.~H.~Conway and R.~K.~Guy, ``Sets of natural numbers with distinct sums'', \emph{Notices Amer.\ Math.\
Soc.}\ 15 (1968), 345. W.~F.~Lunnon, ``Integer sets with distinct subset-sums'', \emph{Math.\ Comp.}\ 50 (1988), 297--320. T.~Bohman, ``A sum packing problem of
Erd\H os and the Conway--Guy sequence'', \emph{Proc.\ Amer.\ Math.\ Soc.}\ 124 (1996), 3627--3636, and ``A construction for sets of integers with distinct subset sums'',
\emph{Electron.\ J.\ Combin.}\ 5 (1998), R3. N.~D.~Elkies, \emph{J.\ Combin.\ Theory Ser.~A} 41 (1986), 89--94. I.~Aliev, ``Siegel's lemma and sum-distinct sets'',
\emph{Discrete Comput.\ Geom.}\ 39 (2008), 59--66. Q.~Dubroff, J.~Fox and M.~W.~Xu, ``A note on the Erd\H os distinct subset sums problem'', \emph{SIAM J.\ Discrete Math.}\
35 (2021), 322--324 (Chapter~\ref{ch:reading}). S.~Steinerberger, ``Some remarks on the Erd\H os distinct subset sums problem'', \emph{Int.\ J.\ Number Theory} (2023), with a history of
the constants.
\item \emph{Data.} OEIS A276661 (values of $a(n)$, Lunnon's computations, Dyson's computation of $a(10)$ in October 2025).
\item \emph{The tools and the community.} K.~Kakaes, ``How AI tore through a mathematical community'', \emph{Quanta Magazine}, 3 August 2026 (the history of
\texttt{erdosproblems.com} since early 2023, the discussion forum of August 2025, the use of language models for literature search and proofs).
On MathOverflow's founding in September 2009: the Simons Foundation article ``The Global Math Commons'' (2011). P.~Ginsparg, ``It was twenty years ago
today\ldots'', arXiv:1108.2700 (2011), on the arXiv's start in August 1991. N.~J.~A.~Sloane, ``The On-Line Encyclopedia of Integer Sequences'', \emph{Notices Amer.\
Math.\ Soc.}\ 65 (2018), 1062--1074, on the 1964 card file.
\item \emph{The second moment method.} N.~Alon and J.~H.~Spencer, \emph{The Probabilistic Method}, 4th ed., Wiley, 2016, Chapter~4, which opens with this
problem.
\end{itemize}
